\section{Optimizations}
\label{sec:optimization}

Our implementation of relational \ematching
 using generic join is simple (under 500 lines),
 but that does not preclude having several optimizations
 important for practical performance.

\subsection{Degenerate Patterns}
\label{sec:degenerate}

Not all patterns correspond to conjunctive queries that involve relational
joins.
Non-nested patterns (whether linear or non-linear) will produce
 relational queries without any joins:
\begin{align*}
  f(\alpha, \beta) &\leftrightarrow R_f(\mathit{root}, \alpha, \beta) \\
  f(\alpha, \alpha) &\leftrightarrow R_f(\mathit{root}, \alpha, \alpha)
\end{align*}
The corresponding query plan is simply a
 scan of a relation with possible filtering.
For these queries,
 generic join (or any other join plan) offers no benefit,
 and building the indices for generic join incurs unnecessary overhead.
A relational \ematching implementation (or any other kind)
 should have a ``fast path''
 for these relatively common types of queries
 that simply scans the \egraph/database for matching \enodes/tuples.
For this reason,
 we exclude these kinds of patterns from our evaluation in \autoref{sec:eval}.

\subsection{Variable Ordering}

Different variable orderings can dramatically affect performance
 for generic join~\cite{eval-wcoj, emptyheaded},
 so choosing a variable ordering is important.
Compared to join plans for binary joins,
 query plans for generic join is much less studied.
In relational \ematching\ we choose a variable ordering
 using two simple heuristics:
First, we prioritize variables that occur in many relations,
 because the intersected set of many relations is likely to be smaller.
Second, we prioritize variables that occur in small relations,
 because intersections involving a small relations are also likely to be smaller.
Performing smaller intersections first can prune down the search space early.

Using these two heuristics,
 the optimizer is able to find more efficient query plans
 than the top-down search of backtracking-based \ematching.
This is even true for linear patterns,
 where our relational \ematching\ has no more information than
 \ematching, but it does have more flexibility.
Consider the linear pattern $f(g(h(\alpha)))$ compiled to
the query $R_f(\mathit{root}, x), R_g(x, y), R_h(y, \alpha)$.
When there are very few $h$-application \enodes\ in the \egraph,
 $R_h$ will be small.
The variable ordering $[y, x, \mathit{root}, \alpha]$
 takes advantage of this by
 intersecting $R_g.y \cap R_h.y$ first,
 resulting in an intersection no larger than $R_h$.
This ``bottom-up'' traversal is not possible in
 conventional \ematching.

%  For example, for queries like $+(a, \alpha)$, our optimizer is able to
% exploit the fact that there is only one atom in $R_a$ (i.e., the relation
% representing constant $a$), so enumerating this relation and use attributes of
% $R_a$ to prune $R_+$ will immediately produces all valid substitutions (assuming
% the indices on $R_+$ is built).
% Meanwhile, in the top-down fashion of backtracking-based \ematching, the
% matching procedure will need to go through all atoms of $R_+$ and check if
% constant $a$ resides in their second child. Such query plan examines much more
% unnecessary atoms in $R_+$, because only a small number of $+$-application
% \enodes{} are connected to $a$.

\subsection{Functional Dependencies}
\label{sec:fd}

Functional dependencies describe the dependencies between columns. For example,
a functional dependency on relation $R(y, x_1, x_2, x_3)$ of the form
$x_1,x_2,x_3\rightarrow y$ indicates that
for each tuple of $R$,  the values of $x_1$, $ x_2$, and $x_3$ uniquely
determines the value $y$. Functional dependencies are ubiquitous in relational
\ematching. In fact, every transformed schema of the form $R_f(\textit{e-class},
\textit{arg}_1,\ldots, \textit{arg}_k)$ has a functional dependency from
$\textit{arg}_1,\ldots,\textit{arg}_k$ to $\textit{e-class}$.
When the variable graph formed by functional dependencies is
 acyclic~\footnote{Cyclicity of functional dependencies is unrelated to cyclicity of the query.},
 we can speed up generic join by ordering the variables to follow the
 topological order of the functional dependency graph.
Every conjunctive query compiled from an \ematching\ pattern
 has acyclic functional dependencies,
 because each dependency goes from the \enode's children to the \enode's parent \eclass.
Relational \ematching~can therefore always choose a variable
 ordering that is consistent with the functional dependency.
Our implementation tries to respect functional dependencies,
 but prioritizes larger intersections more.

As an example, consider conjunctive query~\ref{eqn:cyclic-cq} again. It is
synthesized from pattern $f(g(\alpha), h(\alpha))$ and, assuming each relation
has size $N$, an AGM bound of $O(N^{3/2})$. Suppose however that we pick the
variable ordering $\pi$ to be $[\alpha, x, y, \textit{root}]$. For every
possible value of $\alpha$ chosen, there will be at most one possible value for
$x,y,$ and $\textit{root}$ by functional dependency, which can be immediately
determined. This reduces the run time from $O(N^{3/2})$ to $O(N)$.

\subsection{Batching}
Generic join always processes one variable at a time,
 even if multiple consecutive variables are from the same atom.
We find this strategy to be inefficient in practice,
 as it results in deeper recursion that does little useful work.

Consider the query $Q(x, y, z, w) \gets R(x, y, \alpha), S(\alpha, z, w)$.
The right variable ordering places $\alpha$ at the front,
 since it is the only intersection.
We observe that variables that only appear
 in one atom can be ``batched''
 with others that only appear in the same atom.
Batched variables are treated as a single variable
 in the trie and intersections.
So instead of variable ordering $[\alpha, x, y, z, w]$,
 we can use $[\alpha, (x, y), (z, w)]$.
This lowers the recursion depth of generic join (from 5 to 3)
 and improves data locality by reducing pointer-chasing.

% For example, suppose the algorithm need to loop over a 2-level
%  trie with two nested loops (without intersection),
%  where the outer loop is large but the inner loop is always small.
% Each loop incurs overhead due to caching,
%  processing the trie with nested loops is slower than
%  simply scanning a flattened vector.
% For this reason, we flatten the last levels of each trie where
%  the variables do not participate in any intersection,
%  thereby ``batching'' the variables.

%%% Local Variables:
%%% mode: latex
%%% TeX-master: "main"
%%% End:
